\documentclass[a4paper,twocolumn,11pt]{quantumarticle}
\pdfoutput=1
\usepackage[utf8]{inputenc}
\usepackage[english]{babel}
\usepackage[T1]{fontenc}
\usepackage{amsmath}
\usepackage{amssymb}
\usepackage{amsthm}
\usepackage{hyperref}
\usepackage{graphicx}
\usepackage{booktabs}
\usepackage{multirow}
\usepackage{xcolor}
\usepackage{algorithm}
\usepackage{algpseudocode}

\newtheorem{theorem}{Theorem}
\newtheorem{proposition}[theorem]{Proposition}
\newtheorem{corollary}[theorem]{Corollary}
\newtheorem{lemma}[theorem]{Lemma}
\newtheorem{conjecture}[theorem]{Conjecture}
\newtheorem{definition}[theorem]{Definition}
\newtheorem{remark}[theorem]{Remark}
\newtheorem{observation}[theorem]{Observation}

\begin{document}

\title{Structure-Dependent Classical Simulability of Quantum Algorithms via GF(2) Null Space Analysis}

\author{Author Name}
\affiliation{Manning College of Information and Computer Sciences, University of Massachusetts Amherst, 140 Governors Drive, Amherst, Massachusetts 01003, USA}
\email{email}



\date{\today}

\maketitle

\begin{abstract}

Clifford-augmented matrix product states (CAMPS) simulate quantum circuits by combining stabilizer tracking with matrix product states, using optimization-free disentangling (OFD) when the GF(2) null space structure permits.
We classify 14 circuit families into three scaling classes based on the GF(2) nullity ratio $\nu/t$: OFD-friendly (random circuits, $\nu/t \to 0$), OFD-hostile (QFT, VQE, Grover; $\nu/t \to 1$), and intermediate, whose limiting ratios we trace to a backbone sampling mechanism with a birthday-effect prediction verified to within $4.4\%$.
We prove a bond dimension bound $\chi \le 2^\nu$, show that repeated-block circuits are generically OFD-hostile, characterize the Z-type obstruction for controlled-rotation circuits, and establish an exact rank identity $\mathrm{rank}(z_{\mathrm{QFT}}) = N - 1$: the first such result for a structured circuit family.
Validation against 936 simulations yields Pearson $r = 0.95$ between predicted and empirical OFD success rates, and GF(2) pre-screening eliminates 95--99\% of futile OFD attempts for structured algorithms.
We release the first open-source CAMPS implementation (OFD + OBD) in Julia.
\end{abstract}

%% ============================================================
\section{Introduction}
\label{sec:intro}
%% ============================================================

Classical simulation of quantum circuits is essential for validating quantum hardware, developing quantum algorithms, and understanding the boundary between classical and quantum computational power.
The Gottesman--Knill theorem establishes that circuits composed solely of Clifford gates can be efficiently simulated classically using the stabilizer formalism~\cite{Gottesman1998,Aaronson2004}, but universal quantum computation requires non-Clifford gates such as the T-gate.
Understanding which Clifford+T circuits admit efficient classical simulation and which do not is a central question in quantum complexity theory.

Matrix product states (MPS) provide a powerful framework for classical simulation, representing quantum states as tensor networks with bounded bond dimension $\chi$~\cite{Schollwock2011,White1992,Vidal2003}.
The computational cost scales polynomially with $\chi$, but highly entangled states require exponentially large bond dimensions.
Recently, Clifford-augmented matrix product states (CAMPS) have emerged as a method combining stabilizer tracking with MPS to exploit the structure of Clifford+T circuits~\cite{Liu2024,Qian2024,MelloTDVP2025,Mello2024}.
In the CAMPS framework, a quantum state is decomposed as $|\psi\rangle = C|\phi\rangle$ where $C$ is a Clifford circuit tracked via stabilizer tableaux at cost $O(N^2)$ and $|\phi\rangle$ is an MPS whose bond dimension depends on the non-Clifford content.

The critical algorithmic component of CAMPS is the \emph{disentangling} step: when a T-gate is absorbed into the CAMPS representation, its effect on the MPS component can either be handled cheaply via optimization-free disentangling (OFD)~\cite{Liu2024} or expensively via optimization-based disentangling (OBD)~\cite{Qian2024}.
OFD offers speedups of $10^4$--$10^5\times$ over OBD, but succeeds only when an algebraic condition related to the GF(2) null space of a binary matrix is satisfied.
Liu and Clark~\cite{Liu2024} showed that for \emph{random} Clifford+T circuits with $t \le N$ T-gates on $N$ qubits, OFD succeeds with high probability, establishing that such circuits are efficiently simulable.
However, the circuits practitioners care about, such as QFT, VQE, Grover search, and QAOA, are highly structured, and the behavior of OFD on these circuits has remained an open question.

In this work, we provide a comprehensive answer by studying the GF(2) null space structure across circuit families at scale.
Our contributions are fivefold:

\textbf{(1) GF(2) null space scaling classification.}
We compute the GF(2) matrix for 14 circuit families at system sizes up to $N = 256$ qubits, revealing three distinct scaling classes: OFD-friendly circuits where the nullity ratio $\nu/t \to 0$, OFD-hostile circuits where $\nu/t \to 1$, and intermediate circuits where $\nu/t$ converges to a constant.
This classification, which requires only stabilizer operations and runs in seconds even at $N = 256$, predicts CAMPS simulability without performing the expensive MPS simulation.

\textbf{(2) Structural theorems.}
We formalize the connection between GF(2) rank and CAMPS bond dimension (Theorem~\ref{thm:bond_dim}), prove that repeated-block circuits are OFD-hostile with nullity growing linearly in repetitions (Proposition~\ref{prop:repeated_block}), and characterize the Z-type obstruction that makes controlled-rotation circuits produce GF(2)-trivial rows (Observation~\ref{prop:z_type}).
For QFT, we prove that each controlled rotation contributes $O(1)$ independent GF(2) rows (Proposition~\ref{prop:qft_rotation}), and establish an exact rank identity $\mathrm{rank}(z_{\mathrm{QFT}}) = N - 1$ (Theorem~\ref{thm:qft_rank}).
The proof shows that the sequential structure of the QFT ensures the last qubit never acquires $X/Y$ content in any twisted Pauli, while each of the other $N - 1$ qubits contributes exactly one independent row.
For the intermediate class, we identify a backbone sampling architecture and prove that the GF(2) rank is bounded by the backbone's $G$-block rank (Proposition~\ref{prop:backbone_bound}), with the classical birthday effect predicting $\nu/t \to 1/e$ for full-rank backbones---a prediction matching two families to within $4.4\%$ at $N = 64$, with monotonic convergence toward $1/e$ (Corollary~\ref{cor:birthday}).

\textbf{(3) Practical GF(2)-guided simulation improvements.}
We demonstrate that computing the GF(2) matrix as a pre-screening step eliminates 95--99\% of futile OFD attempts for structured algorithms and 12--18\% for random circuits, providing immediate computational savings.

\textbf{(4) Empirical validation.}
We validate GF(2) predictions against 936 small-scale CAMPS simulations (OFD + OBD), obtaining Pearson $r = 0.95$ between predicted and empirical OFD success rates.
Under pure OFD (no OBD fallback), the structural bound $\chi \le 2^{t-m}$ holds across all 330 tested circuits, with greedy circuit-order OFD achieving the optimal bound $\chi \le 2^\nu$ in 96.7\% of cases.

\textbf{(5) Open-source implementation.}
We provide an open-source implementation of the full CAMPS simulation pipeline (OFD + OBD) in Julia, built on QuantumClifford.jl and ITensorMPS.jl, together with GF(2) analysis tools and benchmarking scripts for all 14 circuit families.

The remainder of this paper is organized as follows.
Section~\ref{sec:background} reviews CAMPS formalism, GF(2) characterization, and related work.
Section~\ref{sec:classification} presents the scaling classification.
Section~\ref{sec:structure} develops the formal results.
Section~\ref{sec:practical} describes practical improvements.
Section~\ref{sec:validation} provides small-scale empirical validation.
Sections~\ref{sec:discussion}--\ref{sec:conclusion} discuss implications and conclude.

%% ============================================================
\section{Background}
\label{sec:background}
%% ============================================================

\subsection{Clifford-Augmented Matrix Product States}

A CAMPS representation decomposes an $N$-qubit quantum state as
\begin{equation}
\label{eq:camps}
|\psi\rangle = C|\phi\rangle,
\end{equation}
where $C$ is a Clifford circuit represented by a stabilizer tableau~\cite{Aaronson2004} and $|\phi\rangle$ is a matrix product state with bond dimension $\chi$.
The Clifford component supports arbitrary stabilizer entanglement at cost $O(N^2)$, while the MPS captures residual non-stabilizer structure.

When a T-gate acts on qubit $q_k$, we use the decomposition $T = \alpha I + \beta Z$ (with $\alpha = e^{i\pi/8}\cos(\pi/8)$, $\beta = -ie^{i\pi/8}\sin(\pi/8)$) to write
\begin{equation}
\label{eq:tgate_absorption}
T_{q_k} C|\phi\rangle = C'(\alpha I + \beta \bar{Z}^{[k]})|\phi\rangle,
\end{equation}
where $C'$ is the updated Clifford circuit and $\bar{Z}^{[k]} = C_k^\dagger Z_{q_k} C_k$ is the \emph{twisted Pauli string}, defined as the image of $Z_{q_k}$ under conjugation by the cumulative Clifford circuit $C_k$ preceding the $k$-th T-gate.
The operator $(\alpha I + \beta \bar{Z}^{[k]})$ must be applied to $|\phi\rangle$; if $\bar{Z}^{[k]}$ has multi-qubit support, this generically increases the bond dimension.

\emph{Optimization-free disentangling (OFD)}~\cite{Liu2024} avoids this bond dimension increase when there exists a qubit $i$ that is simultaneously (a) in a product state within $|\phi\rangle$ (a ``free'' qubit), and (b) has an $X$ or $Y$ component in $\bar{Z}^{[k]}$.
When such a qubit exists, a controlled-Pauli disentangler (a Clifford gate absorbed into the frame $C'$) reduces $\bar{Z}^{[k]}$ to a single-qubit operator, absorbing the T-gate at no bond dimension cost.
When no such qubit exists, OFD fails, and the naive absorption of $(\alpha I + \beta \bar{Z}^{[k]})$ into the MPS at most doubles the bond dimension.
An alternative fallback is \emph{optimization-based disentangling (OBD)}~\cite{Qian2024}, which searches over two-qubit Clifford gates to minimize entanglement growth, but at $10^4$--$10^5\times$ higher computational cost.

More generally, the CAMPS framework applies to any non-Clifford gate expressible as $\alpha I + \beta P$ for a Pauli operator $P$.
An arbitrary $R_z(\theta)$ gate decomposes as $R_z(\theta) = e^{i\theta/2}(\cos(\theta/2)\, I - i\sin(\theta/2)\, Z)$, yielding the same OFD/OBD structure with $\bar{Z}^{[k]} = C_k^\dagger Z_{q_k} C_k$.
Throughout this work, we refer to any non-Clifford $R_z(\theta)$ gate as a ``non-Clifford gate'' in the GF(2) analysis; the term ``T-gate'' is used when the angle is specifically $\pi/4$.

\subsection{GF(2) Characterization of OFD}
\label{sec:gf2_background}

The success or failure of OFD across all non-Clifford gates in a circuit is governed by a binary matrix over GF(2).

\begin{definition}[GF(2) matrix]
\label{def:gf2_matrix}
For a Clifford+non-Clifford circuit on $N$ qubits with $t$ non-Clifford gates, the \emph{GF(2) matrix} $z \in \mathbb{F}_2^{t \times N}$ has entries
\begin{equation}
\label{eq:gf2_entry}
z_{k,j} = \begin{cases} 1 & \text{if qubit $j$ of } \bar{Z}^{[k]} \text{ has Pauli type $X$ or $Y$,} \\ 0 & \text{if qubit $j$ of } \bar{Z}^{[k]} \text{ has Pauli type $I$ or $Z$,} \end{cases}
\end{equation}
where $\bar{Z}^{[k]} = C_k^\dagger Z_{q_k} C_k$ is the $k$-th twisted Pauli string and $C_k$ is the Clifford sub-circuit consisting of all Clifford gates preceding the $k$-th non-Clifford gate.
\end{definition}

The GF(2) rank of $z$ determines the maximum number of non-Clifford gates guaranteed disentanglable via OFD, and the nullity $\nu = t - \mathrm{rank}(z)$ provides an upper bound on the number that must fall back to expensive alternatives.
Liu and Clark~\cite{Liu2024} showed that for random Clifford+T circuits with $t = N$ T-gates at depth $d \ge N/4$, $\mathrm{rank}(z) \approx N - 0.85$ with high probability, where $0.85$ is the expected nullity of a random $N \times N$ binary matrix over $\mathbb{F}_2$.

\subsection{Stabilizer Tableau Computation}
\label{sec:tableau}

The GF(2) matrix can be computed efficiently using stabilizer tableaux~\cite{Aaronson2004}.
A stabilizer tableau for an $N$-qubit Clifford circuit $C$ stores $2N$ generators: destabilizers $\{d_i\}_{i=1}^N$ and stabilizers $\{s_i\}_{i=1}^N$, each an $N$-qubit Pauli operator encoded as a pair of $N$-bit binary vectors $(x_i, z_i)$ plus a phase bit.
The stabilizer row for qubit $q$ encodes the image $C^\dagger Z_q C$, with the $x$-bits giving exactly row $q$ of the GF(2) matrix.
Concretely, $z_{k,j}$ equals the $x$-bit at position $j$ of the $(N+q_k)$-th row of the tableau after applying all Clifford gates preceding the $k$-th non-Clifford gate~\cite{Dehaene2003}.

This connection to the symplectic formulation of the Clifford group is central to our structural analysis.
The Clifford group on $N$ qubits has a faithful representation as the binary symplectic group $\mathrm{Sp}(2N, \mathbb{F}_2)$~\cite{Dehaene2003}.
Each Clifford operator $C$ corresponds to a $2N \times 2N$ symplectic matrix
\begin{equation}
\label{eq:symplectic}
S_C = \begin{pmatrix} A & B \\ G & H \end{pmatrix},
\end{equation}
where the blocks are $N \times N$ binary matrices satisfying the symplectic condition $S_C^T \Omega S_C = \Omega$ with $\Omega = \bigl(\begin{smallmatrix} 0 & I \\ I & 0 \end{smallmatrix}\bigr)$.
The image of $Z_{q}$ under conjugation $C^\dagger Z_q C$ has $x$-bits given by row $q$ of the $G$-block~\cite{Dehaene2003}.
Thus $z_{k,j} = G^{(k)}_{q_k,j}$, where $G^{(k)}$ is the $G$-block of the symplectic matrix for $C_k$.

The linearity of the symplectic representation, in which composition of Cliffords corresponds to matrix multiplication over $\mathbb{F}_2$, is the key algebraic fact underlying our structural results.

\subsection{Related Work}

CAMPS was introduced by Qian, Huang, and Qin~\cite{Qian2024} for ground-state simulation of quantum spin chains, with OBD as the disentangling strategy.
Liu and Clark~\cite{Liu2024} developed the OFD algorithm and the GF(2) characterization for circuit simulation, establishing efficient simulability of random Clifford+T circuits.
Extensions include Clifford-augmented time-dependent variational principle (CA-TDVP)~\cite{Qian2024b}, Clifford-dressed TDVP with iterative two-qubit Clifford sweeps for entanglement reduction during Hamiltonian dynamics~\cite{MelloTDVP2025}, and applications to noisy circuit simulation~\cite{Mello2024}.

Competing classical simulation approaches include stabilizer tensor networks (STN)~\cite{MasotLlima2024}, which decompose states into sums of stabilizer states with MPS-like structure, and low-rank stabilizer decomposition methods~\cite{Bravyi2019,BravyiGosset2016}, which extend the stabilizer formalism to circuits with non-Clifford gates.
Pauli propagation methods~\cite{PauliProp2019} offer alternative approaches via Heisenberg-picture evolution.

Our work connects to the broader study of nonstabilizerness (magic) and its relationship to classical simulability.
Zhang and Zhang~\cite{Zhang2025} established sharp complexity transitions as a function of \emph{magic depth}, the number of non-Clifford layers, showing that T-depth 1 circuits allow polynomial-time Pauli expectation evaluation, while T-depth 2 is GapP-complete.
Fux \emph{et al.}~\cite{Fux2025} analyzed disentangling magic states through a quantum error correction lens, showing complete disentangling when $t \lesssim N - 1.6$.
Frau \emph{et al.}~\cite{Frau2025} classified ground states of conformal field theories into fully, locally, and non-locally Clifford-disentanglable (fLCD, LCD, nLCD) classes based on how effectively local Clifford sweeps reduce entanglement, a state-level classification that has suggestive parallels with our circuit-level GF(2)-based taxonomy.

The nonstabilizerness entanglement entropy framework~\cite{Huang2025,Turkeshi2025} provides complementary measures of how non-Clifford gates generate entanglement in the MPS component, while T-gate compilation theory~\cite{Ross2016,Vandaele2025} governs the T-gate count scaling that determines the dimensions of our GF(2) matrices.

%% ============================================================
\section{GF(2) Null Space Scaling Classification}
\label{sec:classification}
%% ============================================================

We now present our main empirical result: a scaling classification of quantum circuit families based on the asymptotic behavior of the GF(2) nullity ratio $\nu(N)/t(N)$ as the system size $N$ grows.

\subsection{Methodology}

For each circuit family and system size $N$, we construct the gate-level circuit, build the GF(2) matrix $z$ via stabilizer tableau propagation (Definition~\ref{def:gf2_matrix}), and compute its rank via Gaussian elimination over $\mathbb{F}_2$.
Each circuit family uses a standard Clifford+non-Clifford decomposition: random circuits inject T-gates ($R_z(\pi/4)$) at specified positions; structured state preparations (GHZ, Bell, graph, cluster, surface code) use the standard Toffoli decomposition requiring 7 T-gates per Toffoli~\cite{Selinger2013}; and QFT decomposes each controlled-$R_m$ gate ($m \ge 2$) as
\begin{equation}
\label{eq:cr_decomp}
R_z\!\left(\frac{\theta}{2}\right)_t \cdot \mathrm{CNOT}(c,t) \cdot R_z\!\left(\!-\frac{\theta}{2}\right)_t \cdot \mathrm{CNOT}(c,t) \cdot R_z\!\left(\frac{\theta}{2}\right)_c
\end{equation}
(gates listed in temporal order, left to right) with exact angle $\theta = 2\pi/2^m$, yielding three $R_z$ gates (two on the target qubit $t$ and one on the control qubit $c$) and two CNOTs per controlled rotation, where each $R_z$ with angle not a multiple of $\pi/2$ is treated as a single non-Clifford gate in the GF(2) analysis.\footnote{The first $R_z(\theta/2)_t$ on the target qubit occurs before either CNOT, when the target has not yet received its Hadamard, producing a Z-type twisted Pauli (zero $x$-bits) that contributes a zero row to the GF(2) matrix. The effective rank contribution per controlled rotation comes from the remaining two $R_z$ gates (Proposition~\ref{prop:qft_rotation}). The VQE benchmark uses a heuristic Clifford+T decomposition of $R_y(\theta)$ gates that does not precisely implement the intended rotation angles. However, VQE's Class~II membership follows from Proposition~\ref{prop:repeated_block} independently of the specific decomposition: any circuit with $r$ repeated blocks of $g > N$ non-Clifford gates satisfies $\nu/t \ge 1 - 1/r$ regardless of gate angles.}

This analysis relies \emph{only} on stabilizer tableau operations and requires no MPS computations.
The cost is $O(N^2 t)$ for tableau propagation plus $O(\min(t,N)^2 \max(t,N))$ for GF(2) elimination, completing in seconds even for the largest configurations ($N = 256$, $t \approx 80{,}000$ for QFT).
This efficiency is central to our methodology: GF(2) analysis serves as a fast predictor of CAMPS simulability, running orders of magnitude faster than the actual CAMPS simulation.

We benchmark 14 circuit families spanning random circuits, standard quantum algorithms, and structured state preparations:
\emph{Random circuits}: random Clifford+T with all-to-all (A2A) and brick-wall connectivity, with $t = N$ T-gates;
\emph{Quantum algorithms}: QFT, Grover search, VQE with hardware-efficient ansatz, QAOA for MaxCut, Bernstein--Vazirani, Deutsch--Jozsa, Simon's algorithm;
\emph{Structured states}: GHZ, Bell, graph, cluster, and surface code states.
System sizes range from $N = 4$ to $N = 64$ qubits for most families, with QFT and VQE extending to $N = 256$, and 4 random instances per configuration for stochastic families.

\subsection{Three Scaling Classes}

\begin{figure}[t]
    \centering
    \includegraphics[width=\linewidth]{figure1.png}
    \caption{GF(2) null space scaling classification. The nullity ratio $\nu/t$ as a function of qubit count $N$ reveals three distinct classes: OFD-hostile circuits (top band, $\nu/t \to 1$), intermediate circuits, and OFD-friendly circuits (bottom, $\nu/t \to 0$). The separation sharpens with increasing $N$, establishing a predictive classification of CAMPS simulability. QAOA uses angles $\gamma = \pi/4$, $\beta = \pi/8$; of the $5N/2$ total rotation gates, only the $N$ mixer rotations are non-Clifford, but the GF(2) analysis processes all rotation gates, yielding $\nu/t \to 3/5$. QFT includes all controlled-$R_m$ rotations ($m \le N$), giving $\nu/t = 0.997$ at $N = 256$ (Table~\ref{tab:qft_rank}).}
    \label{fig:scaling}
\end{figure}

\begin{figure}[t]
    \centering
    \includegraphics[width=\linewidth]{figure2.png}
    \caption{Absolute nullity $\nu(N)$ on a log--log scale. QFT and VQE exhibit super-linear nullity growth ($\nu \sim N^2$ for QFT), while random circuits maintain constant nullity $\nu \approx 0.85$ independent of $N$. The divergent scaling between families confirms that the classification in Fig.~\ref{fig:scaling} reflects fundamentally different algebraic structure, not merely finite-size effects.}
    \label{fig:nullity_growth}
\end{figure}

Figure~\ref{fig:scaling} displays the nullity ratio $\nu(N)/t(N)$ as a function of system size for all 14 families, revealing a striking three-class structure:

\textbf{Class~I (OFD-friendly):} Random Clifford+T circuits with $t = N$ T-gates.
As $N$ grows, the nullity ratio decreases toward zero, with $\nu/t \approx 0.85/N$, because the absolute nullity remains approximately constant at $\nu \approx 0.85$, corresponding to the expected nullity of a random $N \times N$ binary matrix~\cite{Liu2024}.
These circuits are efficiently simulable via CAMPS with OFD; the bond dimension bound gives $\chi_{\mathrm{OFD}} \le 2^{0.85} < 2$.

\textbf{Class~II (OFD-hostile):} QFT, VQE (multi-layer), and Grover search (multi-iteration).
The nullity ratio converges to 1 as $N$ grows, meaning essentially all non-Clifford gates fall in the null space and OFD provides no benefit.
QFT has $\nu/t > 0.91$ for $N \ge 8$, reaching $0.997$ at $N = 256$ (Table~\ref{tab:qft_rank}); VQE and Grover approach $\nu/t \ge 0.99$.
For these circuits, CAMPS simulation must rely entirely on the slow OBD procedure or alternative strategies.

\textbf{Class~III (intermediate):} QAOA, Bernstein--Vazirani, Deutsch--Jozsa, Simon's, Bell states, graph states, cluster states, and surface codes.
The nullity ratio converges to a constant strictly between 0 and 1 for the families with analytical predictions; surface codes show a slowly increasing ratio ($\nu/t \approx 0.78$ at $N = 50$) whose asymptotic class requires further study.
These circuits derive partial benefit from OFD; the fraction of disentanglable non-Clifford gates is neither negligible nor dominant.
Most Class~III families share a \emph{backbone sampling} architecture that quantitatively explains their limiting nullity ratios (Section~\ref{sec:backbone}); QAOA is structurally distinct, with its limiting ratio $\nu/t \to 3/5$ arising from rank saturation at $N$ with $t = 5N/2$ total rotation gates (Remark~\ref{rem:qaoa}).
GHZ states have backbone rank $r_G = 1$, giving $\mathrm{rank}(z) = 1$ for all $N$ tested and $\nu/t = (N-1)/N \to 1$, placing them at the boundary of Class~II (see Section~\ref{sec:backbone}).

Figure~\ref{fig:nullity_growth} shows the absolute nullity on a log--log scale, revealing the growth rates underlying each class.
QFT exhibits quadratic nullity growth $\nu \sim N^2$ (matching its $O(N^2)$ non-Clifford gate count), while random circuits maintain $\nu \approx 0.85$ independent of $N$.
The intermediate class shows linear to sub-linear growth.
These divergent scaling behaviors confirm that the classification reflects fundamentally different algebraic structure, not finite-size effects.

\begin{definition}[Scaling classes]
\label{def:classes}
A circuit family $\mathcal{F}$ belongs to Class~I if $\lim_{N\to\infty} \nu(N)/t(N) = 0$, Class~II if $\lim_{N\to\infty} \nu(N)/t(N) = 1$, and Class~III if the limit exists and lies in $(0,1)$.
\end{definition}

The classification has immediate practical significance: Class~I circuits are efficiently simulable via CAMPS with OFD, Class~II circuits gain no OFD benefit and require alternative simulation strategies, and Class~III circuits benefit from OFD in proportion to their limiting nullity ratio.
Importantly, the class membership can be determined \emph{before} any MPS simulation, using only the $O(N^2 t)$-cost GF(2) analysis.

\emph{Role of non-Clifford gate density.}
Since $\mathrm{rank}(z) \le N$, any circuit family with $t(N)/N \to \infty$ automatically has $\nu/t \to 1$ and is Class~II.
Our Class~II algorithms, including QFT ($t = 3N(N-1)/2$), multi-layer VQE ($t = O(rN)$ with $r > 1$), and multi-iteration Grover ($t = O(\sqrt{2^N} \cdot N)$), all have non-Clifford gate counts that grow faster than $N$, and therefore their Class~II membership follows directly from the rank bound (Proposition~\ref{prop:repeated_block}).
Random circuits with $t \gg N$ would likewise be Class~II.
The classification's non-trivial content lies in the \emph{intermediate} regime $t = \Theta(N)$: here, random circuits achieve near-full rank $\mathrm{rank} \approx N$ (Class~I), while the algebraic structure of specific algorithms could in principle produce rank-deficient GF(2) matrices even at $t = N$.
The structural theorems in Section~\ref{sec:structure} go beyond the coarse $t/N$ ratio to explain \emph{why} rank saturates (repeated-block structure, Z-type obstruction) and \emph{how} it grows for QFT (Theorem~\ref{thm:qft_rank}), providing finer-grained predictions than the rank bound alone.

%% ============================================================
\section{Structural Analysis}
\label{sec:structure}
%% ============================================================

We now develop formal results explaining the scaling classification.
We first formalize the bond dimension bound implicit in Liu and Clark's framework, then prove a repeated-block result explaining Class~II membership for VQE and Grover, characterize the Z-type obstruction mechanism, prove a per-rotation rank bound for QFT, establish an exact rank identity for the QFT circuit, and identify a backbone sampling mechanism explaining the intermediate class.

\subsection{Bond Dimension Bound}

The following theorem formalizes the bond dimension analysis developed across Liu and Clark's OFD framework~\cite{Liu2024}, connecting the GF(2) nullity directly to worst-case CAMPS bond dimension.

\begin{theorem}[OFD bond dimension bound]
\label{thm:bond_dim}
Let $z \in \mathbb{F}_2^{t \times N}$ be the GF(2) matrix of a Clifford+non-Clifford circuit on $N$ qubits with $t$ non-Clifford gates (Definition~\ref{def:gf2_matrix}).
There exists a gate processing order under which OFD succeeds for at least $\mathrm{rank}_{\mathbb{F}_2}(z)$ gates.
For any processing order in which $m$ gates succeed via OFD, the resulting CAMPS bond dimension satisfies
\begin{equation}
\label{eq:bond_bound}
\chi_{\mathrm{OFD}} \le 2^{t - m}.
\end{equation}
In particular, taking $m \ge \mathrm{rank}(z)$:
\begin{equation}
\label{eq:bond_bound_rank}
\chi_{\mathrm{OFD}} \le 2^{t - \mathrm{rank}_{\mathbb{F}_2}(z)} = 2^{\nu},
\end{equation}
where $\nu = \dim(\ker(z))$ is the GF(2) nullity.
\end{theorem}

\begin{proof}
We establish the three claims in sequence.

\emph{Claim 1: At least $r = \mathrm{rank}(z)$ gates can be disentangled.}
Compute the reduced row echelon form (RREF) of $z$, with pivot positions at columns $j_1 < j_2 < \cdots < j_r$.
Process non-Clifford gates corresponding to pivot rows in order $k_1, k_2, \ldots, k_r$.
For gate $k_i$: by the RREF structure, $z_{k_i, j_i} = 1$ and $z_{k_\ell, j_i} = 0$ for all $\ell < i$.
This means qubit $j_i$ has an $X$ or $Y$ component in the twisted Pauli $\bar{Z}^{[k_i]}$.
Since the initial MPS has all qubits in product states, and OFD for gate $k_\ell$ ($\ell < i$) consumes qubit $j_\ell \ne j_i$ (the OFD disentangler modifies only the Clifford frame, not the MPS tensor structure on other qubits), qubit $j_i$ remains free when gate $k_i$ is processed.
Furthermore, the $x$-bit at position $j_i$ in the twisted Pauli $\bar{Z}^{[k_i]}$ is preserved under prior frame modifications.
The OFD disentangler for gate $k_\ell$ ($\ell < i$) is a product of controlled-Pauli gates with control qubit $j_\ell$; in the Aaronson--Gottesman convention, $\mathrm{CNOT}(j_\ell \to q)$ updates $x$-bits as $G[\cdot,q] \leftarrow G[\cdot,q] \oplus G[\cdot,j_\ell]$, while $\mathrm{CPHASE}$ and single-qubit phase gates leave $x$-bits unchanged.
Thus only columns $q$ with $z_{k_\ell, q} = 1$ are modified.
By the RREF structure, $z_{k_\ell, j_i} = 0$ for $i \ne \ell$ (later pivot columns are zeroed in earlier pivot rows), so column $j_i$ is unmodified by $D_\ell$.
By induction over $\ell = 1, \ldots, i-1$, the $x$-bit $z_{k_i, j_i} = 1$ is preserved in the frame-modified tableau.
Therefore OFD succeeds: the disentangler maps $\bar{Z}^{[k_i]}$ to a single-qubit operator on qubit $j_i$, and the bond dimension is unchanged.

\emph{Claim 2: Each OFD failure at most doubles $\chi$.}
When OFD fails for gate $k$, the operator $(\alpha I + \beta \bar{Z}^{[k]})$ is applied directly to $|\phi\rangle$, yielding $\alpha|\phi\rangle + \beta \bar{Z}^{[k]}|\phi\rangle$.
Since $\bar{Z}^{[k]}$ is a tensor product of single-qubit Paulis, it preserves the MPS structure with identical bond dimension~\cite{Schollwock2011}.
By the triangle inequality for MPS bond dimensions, the sum has bond dimension at most $2\chi$.

\emph{Claim 3: Combined bound.}
Starting from $\chi_0 = 1$ (product state), $m$ OFD successes contribute $\chi \to \chi$ (unchanged) and $t - m$ OFD failures contribute $\chi \to 2\chi$ each, giving $\chi \le 2^{t-m} \cdot 1 = 2^{t-m}$.
Since $m \ge r = \mathrm{rank}(z)$ is achievable by Claim~1, the bound $\chi \le 2^{t-r} = 2^\nu$ follows.
\end{proof}

\begin{remark}[Tightness of the rank bound]
\label{rem:matching}
The bound $\chi \le 2^\nu$ is conservative: the actual maximum number of OFD successes can exceed $\mathrm{rank}(z)$.
An OFD success for gate $k$ requires a free qubit $j$ with $z_{k,j} = 1$; each success consumes a distinct qubit.
The maximum number of successes is therefore the size of a maximum matching in the bipartite graph with edges $\{(k,j) : z_{k,j} = 1\}$, which can exceed $\mathrm{rank}_{\mathbb{F}_2}(z)$ in general.
Nevertheless, the GF(2) rank provides an algebraically computable lower bound that closely tracks empirical OFD success (Pearson $r = 0.95$; see Section~\ref{sec:validation}).
For the bond dimension \emph{upper bound}, the distinction is immaterial: $\chi \le 2^{t-m} \le 2^{t - \mathrm{rank}(z)} = 2^\nu$ for any achievable $m \ge \mathrm{rank}(z)$.
\end{remark}

\begin{figure}[t]
    \centering
    \includegraphics[width=\linewidth]{figure3.png}
    \caption{Bond dimension bound validation (Theorem~\ref{thm:bond_dim}). \textbf{(a)} Log-scale comparison of final bond dimension $\chi$ versus nullity $\nu$, with the bound $\chi \le 2^\nu$ shown as the diagonal. All 936 circuits (OFD + OBD hybrid strategy) satisfy the bound; under pure OFD, the structural bound $\chi \le 2^{t-m}$ holds universally (Section~\ref{sec:bound_tightness}). \textbf{(b)} Tightness ratio $\log_2(\chi)/\nu$ by circuit family, grouped by category (random, algorithm, state preparation). The bound is tightest for random circuits; the looseness for structured algorithms is a finite-size effect (see Section~\ref{sec:bound_tightness}).}
    \label{fig:bound_validation}
\end{figure}

Figure~\ref{fig:bound_validation} validates the bound against 936 CAMPS simulations (Section~\ref{sec:bound_tightness}).

\begin{remark}[Predictive vs.\ descriptive value]
\label{rem:predictive}
The bound's primary value is \emph{predictive}: it enables assessment of CAMPS simulability from the GF(2) matrix alone, before any MPS computation.
In practice, OBD fallback may yield much smaller bond dimensions (Section~\ref{sec:validation}), but the nullity ratio still governs the \emph{time cost} (fraction of gates requiring slow OBD), while the bound governs the worst-case \emph{space cost} (bond dimension).
\end{remark}

\subsection{Repeated-Block Circuits}
\label{sec:repeated_block}

Many quantum algorithms, such as VQE, QAOA, Grover search, and Trotter circuits, are composed of repeated applications of a fixed block.
The following proposition shows that such circuits are generically OFD-hostile whenever the block contains more non-Clifford gates than qubits.

\begin{proposition}[Repeated-block nullity]
\label{prop:repeated_block}
Consider a circuit on $N$ qubits consisting of $r$ repetitions of a block $B$ containing $g$ non-Clifford gates, with Clifford layers $C_1, C_2, \ldots, C_{r-1}$ interleaving the blocks.
The GF(2) rank satisfies
\begin{equation}
\label{eq:rank_upper}
\mathrm{rank}(z) \le N,
\end{equation}
and consequently the nullity satisfies
\begin{equation}
\label{eq:repeated_nullity}
\nu \ge \max(0,\; rg - N).
\end{equation}
When $g \ge N$, the nullity ratio satisfies $\nu/t \ge 1 - 1/r$, which approaches $1$ as $r \to \infty$.
\end{proposition}

\begin{proof}
The GF(2) matrix $z$ has $t = rg$ rows and $N$ columns.
Since $z \in \mathbb{F}_2^{rg \times N}$, the rank satisfies $\mathrm{rank}(z) \le \min(rg, N) \le N$.
Therefore $\nu = rg - \mathrm{rank}(z) \ge rg - N$.

When $g \ge N$: $\nu/t = 1 - \mathrm{rank}(z)/(rg) \ge 1 - N/(rg) \ge 1 - 1/r$.
\end{proof}

While the bound follows from the elementary fact that an $N$-column binary matrix has rank at most $N$, its implications are significant: any circuit with $rg > N$ has OFD-hostile non-Clifford gates \emph{regardless of the specific circuit structure}.
For VQE with $g \approx 2N$ non-Clifford gates per layer and $r = 10$ layers, the nullity ratio satisfies $\nu/t \ge 1 - N/(20N) = 0.95$, explaining VQE's Class~II membership.

The bound can be understood more concretely through the block structure.
The GF(2) matrix decomposes as $z = (z^{(1)\top}, z^{(2)\top}, \ldots, z^{(r)\top})^\top$ where $z^{(k)} \in \mathbb{F}_2^{g \times N}$ collects the rows from block $k$.
Since all rows lie in $\mathbb{F}_2^N$, the combined row space $\mathrm{rowspace}(z) = \sum_{k=1}^r \mathrm{rowspace}(z^{(k)})$ has dimension at most $N$.
When $g \ge N$, the first block's $g$ non-Clifford gates typically produce GF(2) rows spanning all of $\mathbb{F}_2^N$ (i.e., $\mathrm{rank}(z^{(1)}) = N$), in which case $\mathrm{rank}(z) = N$ \emph{independently of $r$}: every non-Clifford gate beyond the first $N$ contributes only to nullity.
Note that the full symplectic representation of each twisted Pauli (both $x$-bits and $z$-bits) transforms linearly under the inter-block Clifford~\cite{Dehaene2003}, but the GF(2) matrix records only the $x$-bit projection, which does not in general transform as a linear image of the previous block's $x$-bits, since symplectic composition mixes $x$-bits and $z$-bits across blocks.

\begin{figure}[t]
    \centering
    \includegraphics[width=\linewidth]{figure4.png}
    \caption{Proposition~\ref{prop:repeated_block} verification. \textbf{(a)} Grover: GF(2) rank saturates at exactly $N$ from $r = 1$ iteration. \textbf{(b)} VQE: identical saturation behavior across all $N$. \textbf{(c)} The ratio $\mathrm{rank}/(g + (r{-}1)N)$ remains below $0.1$, confirming the bound is universally satisfied. The saturation at $\mathrm{rank} = N$ from the first repetition confirms that the single block already spans $\mathbb{F}_2^N$, and subsequent repetitions add zero independent rows.}
    \label{fig:prop2}
\end{figure}

Figure~\ref{fig:prop2} validates this analysis numerically.
For both Grover and VQE at $N = 4, 6, 8, 10, 12$, the GF(2) rank saturates at exactly $N$ from the first repetition ($r = 1$) and remains constant as $r$ increases.
This confirms that within a single block, the $g > N$ non-Clifford gates already span the full $\mathbb{F}_2^N$ row space, and subsequent repetitions add zero independent rows.
The ratio $\mathrm{rank}/[g + (r-1)N]$ (Fig.~\ref{fig:prop2}(c)) remains below $0.1$, showing that the tight bound $\mathrm{rank} \le N$ is substantially stronger than the naive row-counting bound $\mathrm{rank} \le g + (r-1)N$ for these circuits.

\subsection{Z-Type Obstruction}
\label{sec:z_type}

The scaling classification shows \emph{that} structured circuits have high nullity, but does not explain \emph{why} specific non-Clifford gates contribute zero rows to the GF(2) rank.
The following characterization, which is a direct consequence of Definition~\ref{def:gf2_matrix}, identifies the mechanism.

\begin{observation}[Z-type obstruction]
\label{prop:z_type}
The $k$-th non-Clifford gate produces a zero row in $z$ (i.e., $z_{k,j} = 0$ for all $j$) if and only if the twisted Pauli $\bar{Z}^{[k]}$ is Z-type: $\bar{Z}^{[k]} \in \{I, Z\}^{\otimes N}$ (up to global phase).
In the symplectic formulation, this is equivalent to
\begin{equation}
G^{(k)}_{q_k,:} = \mathbf{0},
\end{equation}
where $G^{(k)}$ is the $G$-block of the symplectic matrix of $C_k$.
A Z-type twisted Pauli arises whenever the Clifford circuit $C_k$ preserves the Z-subalgebra on qubit $q_k$: $C_k^\dagger Z_{q_k} C_k \in \{I, Z\}^{\otimes N}$.
\end{observation}

\begin{proof}
By Definition~\ref{def:gf2_matrix}, $z_{k,j} = 0$ for all $j$ if and only if every qubit of $\bar{Z}^{[k]}$ has Pauli type $I$ or $Z$, which is exactly the condition $\bar{Z}^{[k]} \in \{I,Z\}^{\otimes N}$.
By the symplectic representation (Section~\ref{sec:tableau}), the $x$-bits of $C_k^\dagger Z_{q_k} C_k$ are given by row $q_k$ of the $G$-block of $S_{C_k}$.
Hence all $x$-bits vanish if and only if $G^{(k)}_{q_k,:} = \mathbf{0}$.
\end{proof}

This observation explains why circuits built from diagonal gates are OFD-hostile.
All diagonal gates, such as $Z$, $S$, $T$, $R_z(\theta)$, $CZ$, and controlled-$R_z$, commute with $Z$ on each qubit up to phase.
In the symplectic formulation, a diagonal gate has $G = 0$~\cite{Dehaene2003}, meaning it cannot convert a Z-type Pauli to one with $X$ or $Y$ content.
Only gates with a Hadamard component, including $H$ and CNOT, introduce nonzero entries to the $G$-block; in the Heisenberg picture, CNOT maps $Z_{\mathrm{tgt}}$ to $Z_{\mathrm{ctrl}} Z_{\mathrm{tgt}}$, thereby mixing the control's $Z$ into the target.

In QFT, each controlled-$R_m$ gate decomposes via Eq.~\eqref{eq:cr_decomp} into CNOT and $R_z$ gates.
CNOT gates introduce $X/Y$ content to the $G$-block, while diagonal $R_z$ gates do not.
Within a single controlled-$R_m$ decomposition, the two CNOTs bracket the $R_z$ gates, creating a paired structure that constrains the $x$-bit content of the twisted Paulis (see Theorem~\ref{thm:qft_rank} for the precise analysis).

\subsection{QFT Per-Rotation Rank Bound}

The QFT's extreme nullity ratio admits a structural explanation.
We first establish a per-rotation bound, then prove the exact rank identity.

\begin{proposition}[Per-rotation rank contribution in QFT]
\label{prop:qft_rotation}
In the $N$-qubit QFT with each controlled-$R_m$ ($m \ge 2$) decomposed via Eq.~\eqref{eq:cr_decomp} using exact angles, each controlled rotation contributes exactly $3$ non-Clifford $R_z$ gates but at most $O(1)$ linearly independent rows to the GF(2) matrix $z$.
Consequently,
\begin{equation}
\label{eq:qft_per_rotation}
\mathrm{rank}(z_{\mathrm{QFT}}) \le c \cdot \frac{N(N-1)}{2}
\end{equation}
for a constant $c$ independent of $N$, while $t = 3N(N-1)/2$.\footnote{Each of the $N(N-1)/2$ controlled rotations contributes 3 non-Clifford $R_z$ gates via the decomposition in Eq.~\eqref{eq:cr_decomp}: two on the target qubit and one on the control qubit. For $N \ge 35$, the implementation's finite-precision Clifford angle test (tolerance $10^{-10}$) misclassifies rotations with $m \gtrsim 35$ as Clifford due to floating-point underflow, reducing the effective $t$ below $3N(N-1)/2$. The values in Table~\ref{tab:qft_rank} reflect the implementation's effective counts. The rank identity $\mathrm{rank} = N - 1$ holds regardless, since these additional gates cannot increase the rank beyond $N - 1$ (Theorem~\ref{thm:qft_rank}).}
Since both the rank bound and $t$ are $O(N^2)$, Proposition~\ref{prop:qft_rotation} alone gives $\nu/t \ge 1 - c/3 \ge 1/3$.
The stronger conclusion $\nu/t \to 1$ follows from the exact rank identity $\mathrm{rank} = N - 1 = O(N)$ established in Theorem~\ref{thm:qft_rank} below.
\end{proposition}

\begin{proof}
Consider a single controlled-$R_m$ decomposition (Eq.~\eqref{eq:cr_decomp}) in the QFT with control qubit $c = j$ and target qubit $k = j + m - 1$.
The decomposition produces three non-Clifford $R_z$ gates: $R_z(\theta/2)$ on target $k$, then $R_z(-\theta/2)$ on target $k$ (between the two $\mathrm{CNOT}(j,k)$ gates), then $R_z(\theta/2)$ on control $j$.

The first gate $R_z(\theta/2)_k$ acts on the target qubit $k > j$, which has not yet received its Hadamard (Hadamards are applied in order $j = 1, 2, \ldots, N$).
By Observation~\ref{prop:z_type}, this gate produces a Z-type twisted Pauli with all-zero $x$-bits, contributing a zero row to $z$.

The second gate $R_z(-\theta/2)_k$ occurs between the two $\mathrm{CNOT}(j,k)$ gates.
The first $\mathrm{CNOT}(j,k)$ adds column $j$ to column $k$ of the $G$-block, potentially modifying entries in column $k$.
This gate reads row $k$ of the $G$-block and contributes at most one $x$-bit vector.

The third gate $R_z(\theta/2)_j$ acts on the control qubit $j$, which has already received its Hadamard, so the $G$-block has $G[j,j] = 1$.
After the second CNOT restores the $G$-block to its pre-pair configuration, the twisted Pauli has $x$-bits determined by row $j$ of $G$, and thus contributes at most one non-zero $x$-bit vector.

The second and third gates produce $x$-bit vectors spanning a subspace of dimension at most 2 over $\mathbb{F}_2$.
With $N(N-1)/2$ controlled rotations, the bound~\eqref{eq:qft_per_rotation} follows with the constant $c \le 2$.
\end{proof}

\subsection{Exact QFT Rank Identity}

Proposition~\ref{prop:qft_rotation} shows that the QFT rank grows at most as $O(N^2)$, corresponding to the number of rotations rather than the total number of non-Clifford gates.
However, the empirical rank grows much slower: it is exactly $N - 1$.

\begin{figure}[t]
    \centering
    \includegraphics[width=\linewidth]{figure5.png}
    \caption{QFT GF(2) rank scaling (extended to $N = 256$). \textbf{(a)} Despite quadratic non-Clifford gate count $t \sim 3N(N-1)/2$ (orange squares; see footnote in Proposition~\ref{prop:qft_rotation} for finite-precision effects at large $N$), the GF(2) rank satisfies $\mathrm{rank} = N - 1$ exactly (green dashed line, zero residuals) for all $N$ tested. The power-law fit $0.77 \cdot N^{1.06}$ and log-corrected fit are shown for comparison; both deviate systematically. The shaded region between rank and $t(N)$ represents the nullity. \textbf{(b)} Residuals on log scale: the $N - 1$ model has zero residuals (green triangles), while competing models show systematic deviations.}
    \label{fig:qft_scaling}
\end{figure}

\begin{figure}[t]
    \centering
    \includegraphics[width=\linewidth]{figure5b.png}
    \caption{Effective power-law exponent $\alpha_{\mathrm{eff}}$ between consecutive $N$ values, converging monotonically from $1.22$ at $N = 8$ to $1.006$ at $N = 256$. The convergence to $\alpha = 1$ (dashed green line) confirms the exact linear scaling $\mathrm{rank} = N - 1$.}
    \label{fig:qft_exponent}
\end{figure}

\begin{theorem}[QFT GF(2) rank]
\label{thm:qft_rank}
For the $N$-qubit QFT with each controlled-$R_m$ ($m \ge 2$) decomposed via Eq.~\eqref{eq:cr_decomp} using exact angles $\theta = 2\pi/2^m$, the GF(2) matrix satisfies
\begin{equation}
\label{eq:qft_rank}
\mathrm{rank}(z_{\mathrm{QFT}}) = N - 1.
\end{equation}
\end{theorem}

\begin{proof}
The QFT circuit processes qubits sequentially: for each $j = 1, \ldots, N$, a Hadamard is applied to qubit $j$, followed by controlled-$R_m$ gates with control qubit $c = j$ and target qubit $k = j + m - 1$ for $m = 2, \ldots, N - j + 1$.
Each controlled-$R_m$ contributes three non-Clifford $R_z$ gates via Eq.~\eqref{eq:cr_decomp}: two on target qubit $k$ and one on control qubit $j$.
The final iteration $j = N$ applies only a Hadamard (no controlled rotations), producing no non-Clifford gates.

We write $\tilde{P}_k = C_k^\dagger Z_{q_k} C_k$ for the twisted Pauli at the $k$-th non-Clifford gate on qubit $q_k$, and $x^{(k)} \in \mathbb{F}_2^N$ for its $x$-bit vector.

\medskip
\noindent\emph{Upper bound: $\mathrm{rank} \le N - 1$.}\quad
We show $x^{(k)}_N = 0$ for every non-Clifford gate $k$, so all rows lie in the subspace $\{v \in \mathbb{F}_2^N : v_N = 0\}$.

All non-Clifford gates occur during iterations $j = 1, \ldots, N - 1$ (iteration $j = N$ has no controlled rotations).
Before the Hadamard on qubit $N$ at the start of iteration $j = N$, qubit $N$ has never been the target of a Hadamard or $S$ gate.
It participates only through $\mathrm{CNOT}(j, N)$ gates (as target) within the controlled-$R_m$ decompositions where the target qubit is $N$ (i.e., controlled rotations from iteration $j$ with $k = N$), and through $R_z$ gates on qubit $N$ within those same decompositions.

The CNOT update rule for the $G$-block (the $x$-bit sub-matrix of the stabilizer tableau) is: $\mathrm{CNOT}(c, t)$ adds column $c$ to column $t$, i.e., $G[\cdot, t] \leftarrow G[\cdot, t] \oplus G[\cdot, c]$, following the Aaronson--Gottesman convention~\cite{Aaronson2004}.
The only other Clifford gate affecting the $G$-block is the Hadamard $H_q$, which swaps the $x$-bit and $z$-bit columns at position $q$; starting from $G = 0$, the Hadamard on qubit $q$ sets $G[\cdot, q]$ to the former $z$-bit column at position $q$ (which has $G[q, q] = 1$ from the initial $Z_q$ stabilizer, while $G[q', q] = 0$ for $q' \ne q$).

We show $x^{(k)}_N = 0$ for every non-Clifford gate $k$, where $x^{(k)}_N = G^{(k)}[q_k, N]$ (row $q_k$, column $N$).

The only operations that can set entries in column $N$ of $G$ are: (i) $H_N$, which swaps the $x$/$z$ columns at position $N$, and (ii) any $\mathrm{CNOT}(c, N)$ with target $N$, which adds column $c$ to column $N$.
The Hadamard $H_N$ occurs only at iteration $j = N$, after all non-Clifford gates, so it cannot affect any twisted Pauli.
The $\mathrm{CNOT}(j, N)$ gates (within controlled-$R_m$ decompositions targeting qubit $N$) do temporarily modify column $N$: each adds column $j$ to column $N$, setting $G[j, N] = 1$ if $G[j, j] = 1$ (which holds when qubit $j$ has received its Hadamard).
However, these CNOTs come in pairs within each decomposition (Eq.~\eqref{eq:cr_decomp}), and the second CNOT restores column $N$ to its pre-pair state.
The three non-Clifford $R_z$ gates within each decomposition see column $N$ as follows:

\emph{Gate 1} ($R_z(\theta/2)$ on target $N$, before the CNOT pair): column $N$ has not been modified, so $G[N, N] = 0$ and $x_N = 0$.

\emph{Gate 2} ($R_z(-\theta/2)$ on target $N$, between the paired CNOTs): the first $\mathrm{CNOT}(j, N)$ has added column $j$ to column $N$.
This gate reads row $N$ of the $G$-block: $x^{(k)}_N = G[N, N]$.
Since the column operation adds $G[N, j]$ to $G[N, N]$, and $G[N, j] = 0$ (qubit $N$ has not had its Hadamard, so row $N$ of $G$ is all zeros before the CNOT), $G[N, N]$ remains $0$, giving $x_N = 0$.

\emph{Gate 3} ($R_z(\theta/2)$ on control $j$): this occurs after the second CNOT restores column $N$.
This gate reads row $j$ of the $G$-block: $x^{(k)}_N = G[j, N] = 0$ (restored by the CNOT pair).

For non-Clifford gates on qubits $q_k \ne N$ from decompositions \emph{not} targeting qubit $N$ (i.e., $\mathrm{CNOT}(j, k)$ with $k < N$): these CNOTs do not modify column $N$ at all, so $G[q_k, N]$ can only be nonzero if a prior unpaired CNOT targeting $N$ left a residue.
Since all $\mathrm{CNOT}(j, N)$ gates occur in cancelling pairs and column $N$ is restored after each pair, no residue accumulates.

Therefore $x^{(k)}_N = 0$ for all non-Clifford gates $k$, and $\mathrm{rank} \le N - 1$.

\medskip
\noindent\emph{Lower bound: $\mathrm{rank} \ge N - 1$.}\quad
At the start of iteration $j$ ($1 \le j \le N - 1$), the Hadamard $H_j$ sets $G[j, j] = 1$ (and no other entry in row $j$ of $G$ is nonzero at this point, since prior CNOT pairs have restored the $G$-block after each controlled rotation).
Consider the first controlled rotation in iteration $j$: the controlled-$R_2$ (controlled-$S$) with control $c = j$ and target $k = j + 1$.
Its decomposition (Eq.~\eqref{eq:cr_decomp}) produces three non-Clifford gates.
The \emph{third} of these is $R_z(\theta/2)$ acting on control qubit $j$. Since it occurs after the second $\mathrm{CNOT}(j, j{+}1)$ restores the $G$-block, it produces a twisted Pauli $C^\dagger Z_j C$ whose $x$-bit vector is $e_j$, the unit vector with a 1 only in position $j$.
This holds because: (i) the paired $\mathrm{CNOT}(j, j{+}1)$ gates within the decomposition cancel their column modifications (the first adds column $j$ to column $j{+}1$, the second restores it), so after both CNOTs the $G$-block returns to its pre-pair state; (ii) the third $R_z$ gate acts on control qubit $j$ and reads row $j$ of the $G$-block, which has $G[j, j] = 1$ from $H_j$ and $G[j, j'] = 0$ for $j' \ne j$, yielding $x$-bit vector $e_j$.
Note that the first two non-Clifford gates (both on target $k = j + 1$) produce all-zero $x$-bit rows: gate~1 sees row $k$ which is all zeros (qubit $k$ has not yet had its Hadamard), and gate~2 also sees row $k$ which remains all zeros (the $\mathrm{CNOT}(j,k)$ modifies column $k$ by adding column $j$, but this only changes $G[j, k]$, not $G[k, k]$, since $G[j, j] = 1$ is the only nonzero entry in column $j$).

The vectors $e_1, e_2, \ldots, e_{N-1}$ are linearly independent over $\mathbb{F}_2$, so $\mathrm{rank} \ge N - 1$.
Subsequent non-Clifford gates within and after iteration $j$ cannot eliminate these pivot rows, as each $e_j$ occupies a distinct column.
Numerical verification confirms this structure for all $N$ from 4 to 256: the RREF of $z_{\mathrm{QFT}}$ has pivots at columns $1, 2, \ldots, N-1$ with no pivot at column $N$ (Table~\ref{tab:qft_rank}).
\end{proof}

\begin{table}[t]
\centering
\small
\caption{QFT GF(2) rank confirming $\mathrm{rank} = N - 1$ for all $N$ tested. The $t$ column reports the effective non-Clifford gate count as computed by the implementation; for $N \ge 35$, this is slightly below $3N(N{-}1)/2$ because the implementation's Clifford-angle test absorbs rotations $R_m$ with $m \ge 35$ (see Proposition~\ref{prop:qft_rotation} footnote). The effective exponent $\alpha_{\mathrm{eff}}$ between consecutive $N$ values converges to $1$.}
\label{tab:qft_rank}
\begin{tabular}{@{}rrrrrc@{}}
\toprule
$N$ & $t$ & rank & $N{-}1$ & $\nu/t$ & $\alpha_{\mathrm{eff}}$ \\
\midrule
4 & 18 & 3 & 3 & 0.833 & --- \\
8 & 84 & 7 & 7 & 0.917 & 1.22 \\
16 & 360 & 15 & 15 & 0.958 & 1.10 \\
32 & 1{,}488 & 31 & 31 & 0.979 & 1.05 \\
64 & 4{,}656 & 63 & 63 & 0.986 & 1.02 \\
128 & 17{,}424 & 127 & 127 & 0.993 & 1.01 \\
256 & 79{,}824 & 255 & 255 & 0.997 & 1.006 \\
\bottomrule
\end{tabular}
\end{table}

Theorem~\ref{thm:qft_rank} is the first exact GF(2) rank result for a specific structured circuit family.
Liu and Clark~\cite{Liu2024} established that random Clifford+T circuits with $t = N$ have $\mathrm{rank} \approx N - 0.85$; our result shows that the QFT, despite having $t = \Theta(N^2)$ non-Clifford gates, achieves only rank $N - 1$---just one dimension short of the maximum possible rank.
The single missing dimension corresponds to qubit $N$ (the last qubit processed), whose Hadamard gate occurs after all non-Clifford gates, preventing it from contributing $X/Y$ content to any twisted Pauli.

\emph{Connection to magic depth.}
Zhang and Zhang~\cite{Zhang2025} established that computational complexity is governed by magic depth, the number of non-Clifford layers. In particular, they identified a sharp transition: T-depth 1 circuits are polynomial-time simulable, whereas T-depth 2 circuits are GapP-complete.
QFT has magic depth $\Theta(N)$ (non-Clifford gates from $N - 1$ Hadamard-interleaved stages), placing it firmly in the computationally hard regime.
Our GF(2) analysis provides a finer characterization: within the hard regime, the QFT has rank exactly $N - 1$, giving nullity $\nu = t - N + 1 = \Theta(N^2)$ and a bond dimension bound $\chi \le 2^{\Theta(N^2)}$.

\begin{figure}[t]
    \centering
    \includegraphics[width=\linewidth]{figure6.png}
    \caption{QFT incremental rank growth per controlled rotation, shown as cumulative GF(2) rank versus fractional gate position. At $N = 8$ (84 non-Clifford gates), $N = 16$ (360 non-Clifford gates), and $N = 32$ (1{,}488 non-Clifford gates), each controlled rotation contributes a staircase step of $O(1)$ rank despite containing 3 non-Clifford gates. Later rotations contribute diminishing rank increments as the row space approaches saturation at $N - 1$.}
    \label{fig:qft_staircase}
\end{figure}

\subsection{Intermediate Circuits: Backbone Sampling}
\label{sec:backbone}

The preceding subsections provide analytical explanations for Class~I (rank saturation from random sampling) and Class~II (repeated-block structure, Z-type obstruction).
We now address the intermediate class.

\begin{remark}[Backbone sampling architecture]
\label{rem:backbone}
All Class~III families except QAOA share a common structure: a deterministic Clifford backbone (Hadamard and CNOT layers) followed by $t = N$ non-Clifford gates placed on qubits drawn from $\{1, \ldots, N\}$.
The backbone's $G$-block, which is the submatrix of the stabilizer tableau recording the $X$-component transformation, determines which qubits produce nonzero $\mathrm{GF}(2)$ rows when targeted by a T-gate.
Specifically, a T-gate on qubit $q$ produces a GF(2) row equal to row $q$ of the backbone's $G$-block at the time of application.
This contrasts with Class~I (random backbone yielding near-full rank) and Class~II (deterministic backbone with $t \gg N$ or diagonal structure yielding near-zero rank).
\end{remark}

\begin{proposition}[Backbone rank bound]
\label{prop:backbone_bound}
Let a circuit consist of a Clifford backbone with $G$-block $G_{\mathrm{bb}} \in \mathbb{F}_2^{N \times N}$, followed by $t$ non-Clifford gates with no intervening Clifford operations.
Then $\mathrm{rank}(z) \le \mathrm{rank}(G_{\mathrm{bb}})$.
\end{proposition}

\begin{proof}
Each non-Clifford gate on qubit $q$ produces a GF(2) row equal to row $q$ of $G_{\mathrm{bb}}$.
Since all rows of $z$ are drawn from the rows of $G_{\mathrm{bb}}$, the row space of $z$ is contained in the row space of $G_{\mathrm{bb}}$, giving $\mathrm{rank}(z) \le \mathrm{rank}(G_{\mathrm{bb}})$.
\end{proof}

\begin{corollary}[Birthday-effect nullity]
\label{cor:birthday}
If $G_{\mathrm{bb}}$ has $N$ linearly independent rows and $t = N$ non-Clifford gates target qubits chosen uniformly and independently from $\{1, \ldots, N\}$, then the expected number of distinct qubits targeted is $N(1 - (1 - 1/N)^N) \to N(1 - 1/e)$ as $N \to \infty$, giving $\mathrm{rank}(z) \to N(1 - 1/e)$ and $\nu/t \to 1/e \approx 0.368$.
\end{corollary}

\begin{proof}
This is the classical birthday/occupancy result.
The probability that a given qubit is untargeted is $(1 - 1/N)^N \to e^{-1}$. It follows that the expected number of distinct targeted qubits is $N(1 - 1/e)$; because the targeted rows of $G_{\mathrm{bb}}$ are independent, this also equals the expected rank.
The nullity ratio is $\nu/t = (t - \mathrm{rank})/t \to 1 - (1 - 1/e) = 1/e$.
\end{proof}

\emph{Empirical verification.}
The prediction $\nu/t \to 1/e$ is confirmed by the data at $N = 64$:
cluster states achieve $\nu/t = 0.352$ (within $4.4\%$ of $1/e$), graph states $0.352$ (within $4.4\%$), and Bernstein--Vazirani $0.398$ (within $8.2\%$, with slower convergence likely due to the structure of its oracle); the convergence toward $1/e \approx 0.368$ is monotonic with increasing $N$.
Bell states ($\nu/t \approx 0.55$) and Simon's algorithm ($\nu/t \approx 0.61$) exhibit partial backbones, since Hadamard gates act on only $N/2$ qubits. Consequently, the effective backbone rank satisfies $r_{\mathrm{eff}} < N$, leading to higher nullity ratios consistent with a reduced sampling pool.
GHZ states have backbone rank $r_G = 1$ (only the first qubit receives a Hadamard), so $\mathrm{rank}(z) = 1$ for all $N$ tested ($N = 4$ to $64$, 48 data points), giving $\nu/t = (N - 1)/N \to 1$.
This places GHZ at the boundary of Class~II despite having only $t = N$ non-Clifford gates.

\begin{remark}[QAOA structure]
\label{rem:qaoa}
QAOA does not fit the backbone sampling model but has a clean structural explanation.
With angles $\gamma = \pi/4$, $\beta = \pi/8$ (following Farhi et al.~\cite{farhi2014quantum}), the cost-layer rotations $R_z(2\gamma) = R_z(\pi/2) = S$ are Clifford, while the $N$ mixer rotations $R_x(2\beta) = R_x(\pi/4)$ are non-Clifford.
The GF(2) analysis processes all $t = 5N/2$ rotation gates (including the $3N/2$ Clifford cost-layer gates, which contribute zero-rank rows), and the rank saturates at the qubit count: $\mathrm{rank}(z) \approx N$.
Since $t = 5N/2$, the nullity ratio converges to $\nu/t = 1 - N/(5N/2) = 3/5 = 0.6$, consistent across all tested system sizes ($N = 4$ to $64$) with negligible variance from the random 3-regular graph topology.
An implementation that absorbs Clifford-angle rotations into the tableau would reduce the effective gate count to $t = N$; we process all rotation gates uniformly for consistency across families.
This contrasts with backbone sampling families where the birthday effect governs the limiting ratio.
\end{remark}

\begin{remark}[Caveats]
\label{rem:caveats}
Proposition~\ref{prop:backbone_bound} applies strictly to circuits where all non-Clifford gates follow the Clifford backbone without intervening Clifford operations.
In practice, OFD frame updates modify the stabilizer tableau after each disentangling step, and some circuits (e.g., Deutsch--Jozsa) interleave Hadamard gates between T-gates, violating the static backbone assumption.
The birthday-effect prediction (Corollary~\ref{cor:birthday}) is therefore a heuristic for circuits with intervening Cliffords, though it remains empirically accurate for the families tested.
\end{remark}

\begin{conjecture}[Class~III limiting nullity]
\label{conj:class3}
For backbone sampling circuits with $N$ linearly independent productive qubits and $t = N$ uniformly random non-Clifford gate placements, $\nu/t \to 1/e$.
For partial backbones with $r_{\mathrm{eff}} < N$ productive qubits, the limiting nullity ratio exceeds $1/e$, with the exact value determined by the backbone structure.
\end{conjecture}

%% ============================================================
\section{Practical GF(2)-Guided Improvements}
\label{sec:practical}
%% ============================================================

The GF(2) analysis naturally suggests practical improvements to CAMPS simulation.
Since the GF(2) matrix is cheap to compute, it serves as a preprocessing step.

\subsection{Pre-Screening of OFD Attempts}

In standard CAMPS, each non-Clifford gate triggers an OFD attempt: the algorithm checks whether a suitable free qubit exists, constructs the disentangler if possible, and falls back to OBD otherwise.
For OFD-hostile circuits, the vast majority of these attempts fail, wasting the computational cost of the OFD check before the (still necessary) OBD invocation.

GF(2) pre-screening eliminates this waste.
By computing the row echelon form of $z$ before simulation, we identify which non-Clifford gates are guaranteed to succeed with OFD (pivot rows) and which cannot succeed under the RREF ordering (null-space rows).
During simulation, only pivot-row gates invoke OFD; null-space gates proceed directly to OBD (or naive absorption), skipping the futile OFD check entirely.

\begin{figure}[t]
    \centering
    \includegraphics[width=\linewidth]{figure9.png}
    \caption{Fraction of OFD attempts eliminated by GF(2) pre-screening. For OFD-hostile circuits (VQE, Grover, QFT), pre-screening skips 95--99\% of futile OFD checks. Even OFD-friendly random circuits benefit modestly (12--18\%), since pre-screening identifies the small number of null-space gates that would fail OFD regardless.}
    \label{fig:prescreening}
\end{figure}

Figure~\ref{fig:prescreening} quantifies the savings.
For VQE, pre-screening eliminates 99\% of OFD attempts; for Grover, 98\%; for QFT, 95\%.
Even for OFD-friendly random circuits, pre-screening provides a 12--18\% reduction.
The GHZ family sees 83\% elimination, consistent with its 23\% empirical OFD success rate (Table~\ref{tab:ofd_rates} in Section~\ref{sec:validation}).
The computational cost of pre-screening is negligible compared to CAMPS simulation: computing $z$ and its row echelon form takes milliseconds, versus hours for the full simulation.

\subsection{Optimal Qubit Assignment}

Beyond pre-screening, the GF(2) analysis informs the choice of \emph{which free qubit} to use as the OFD control for each non-Clifford gate.
Standard greedy OFD selects the first available free qubit with a nonzero $x$-bit; however, this can consume a qubit needed by a later gate, reducing the total number of OFD successes below $\mathrm{rank}(z)$.
Theorem~\ref{thm:bond_dim} guarantees that an assignment achieving $m \ge \mathrm{rank}(z)$ exists (via the RREF pivot columns), but this optimal assignment requires processing gates in the RREF-determined order, which may differ from circuit order.
Since non-Clifford gates must be processed in circuit order (each twisted Pauli depends on the cumulative Clifford state, including prior OFD frame modifications), the RREF assignment is not directly implementable.
In practice, greedy circuit-order OFD achieves the optimal bound $m \ge \mathrm{rank}(z)$ in 96.7\% of tested circuits, with a rank--$m$ deficit of at most 1--2 in the remaining cases (Section~\ref{sec:bound_tightness}).

%% ============================================================
\section{Small-Scale Empirical Validation}
\label{sec:validation}
%% ============================================================

We validate the GF(2) predictions against full CAMPS simulations on 936 circuits at $N = 4$ to $16$ qubits, providing ground truth for the theoretical framework.

\subsection{Implementation}

We implement CAMPS in Julia using QuantumClifford.jl~\cite{Krastanov2022} for stabilizer tableau operations and ITensorMPS.jl~\cite{Fishman2022} for MPS operations.
We employ a hybrid disentangling strategy: for each non-Clifford gate, OFD is attempted first following Liu and Clark's criterion~\cite{Liu2024}; on failure, we fall back to OBD following Qian \emph{et al.}~\cite{Qian2024}, searching over the $|\mathrm{Cl}_2| = 11{,}520$ two-qubit Clifford operators to minimize the second R\'enyi entropy, using their efficient $O(\chi^3 + |\mathrm{Cl}_2|)$ evaluation method.
This hybrid OFD$\to$OBD pipeline is our implementation choice; Liu and Clark's original algorithm applies the gate directly to the MPS when OFD fails, accepting the bond dimension increase.
The implementation spans approximately 9{,}000 lines of Julia code.

For each circuit, we record the OFD success/failure at every non-Clifford gate, wall-clock simulation time, final MPS bond dimension, and GF(2) rank computed independently.

\subsection{GF(2) Prediction Accuracy}

\begin{table}[t]
\centering
\small
\caption{Summary statistics across 14 circuit families at $N = 4$--$16$ qubits.
OFD\% is the fraction of non-Clifford gates successfully disentangled via OFD; Rank/T is the GF(2)-predicted fraction; the Pearson correlation across all 904 individual circuits (excluding pure-Clifford) is $r = 0.95$.}
\label{tab:ofd_rates}
\begin{tabular}{@{}lrrrrr@{}}
\toprule
Family & $t$ & Rank & $\nu$ & OFD\% & $\chi$ \\
\midrule
Random (A2A) & 12 & 9.6 & 2.4 & 86 & 38.0 \\
Random (Brick) & 12 & 9.4 & 2.6 & 82 & 38.0 \\
Bernstein--Vaz. & 12 & 7.3 & 4.7 & 68 & 2.0 \\
Graph State & 12 & 7.4 & 4.6 & 66 & 2.3 \\
Cluster State & 12 & 7.2 & 4.8 & 65 & 2.1 \\
Bell State & 12 & 5.1 & 6.9 & 56 & 2.4 \\
Deutsch--Jozsa & 12 & 5.8 & 6.8 & 51 & 1.9 \\
Simon's & 12 & 5.1 & 6.9 & 49 & 2.6 \\
Surface Code & 9 & 2.1 & 7.3 & 43 & 1.5 \\
QAOA & 30 & 11.6 & 18.4 & 39 & 99.0 \\
GHZ & 12 & 1.0 & 11.0 & 23 & 1.8 \\
QFT & 44 & 5.0 & 38.7 & 12 & 1.0 \\
Grover & 98 & 4.0 & 94.0 & 2 & 4.0 \\
VQE & 598 & 5.0 & 593.0 & 1 & 6.0 \\
\bottomrule
\end{tabular}
\vspace{1mm}

\footnotesize{Values are averages over circuits with $t > 0$ (excluding pure-Clifford instances) from the 936-circuit benchmark. OFD\% $= \overline{m_i/t_i}$, the per-circuit fraction of non-Clifford gates disentangled by OFD, averaged over the family. Phase~1 families (Random through GHZ) use $N = 8, 12, 16$ with $t = 12$ non-Clifford gates each; Deutsch--Jozsa generates some pure-Clifford circuits (32 of 72) that are excluded. Phase~2 families span varying qubit counts and $t$-counts. $\chi$ = final bond dimension from full CAMPS simulation (OFD + OBD). For QFT, the exact-angle decomposition (Eq.~\eqref{eq:cr_decomp}) at $N = 4$--$8$ yields 3 non-Clifford $R_z$ gates per controlled-$R_k$; the GF(2) rank is exactly $N - 1$ (Theorem~\ref{thm:qft_rank}), and the hybrid OFD+OBD strategy achieves $\chi = 1$ at all tested sizes, confirming that the Clifford frame captures QFT's entanglement structure exactly. QAOA uses fixed angles $\gamma = \pi/4$, $\beta = \pi/8$ at $N = 8, 12, 16$; all $t = 5N/2$ rotation gates (including $3N/2$ Clifford-angle cost-layer $R_z(\pi/2)$) are processed uniformly, yielding $t = 20, 30, 40$ per size (Remark~\ref{rem:qaoa}).
Grover data is from $N = 4$ only. The bound $\chi \le 2^\nu$ (Theorem~\ref{thm:bond_dim}) applies per-circuit, not to averaged quantities.}
\end{table}

Table~\ref{tab:ofd_rates} summarizes the empirical results.
The GF(2) rank ratio $\mathrm{rank}/t$ closely tracks the empirical fraction of non-Clifford gates disentangled via OFD across all families, from 86\% for random all-to-all circuits down to 1\% for VQE.
The Pearson correlation between these quantities, computed at the individual circuit level across all 904 non-trivial circuits, is $r = 0.95$, confirming that the GF(2) matrix provides a precise predictor of OFD behavior.
At the per-family level (14 families, averaging over system sizes), the correlation is $r = 0.98$.
For QFT specifically, the GF(2) rank is exactly $N - 1$ at every system size tested (Theorem~\ref{thm:qft_rank}), and the hybrid OFD+OBD benchmark achieves $\chi = 1$ for all 72 QFT circuits ($N = 4$--$8$), confirming that the Clifford frame perfectly captures QFT's entanglement.
In the hybrid simulation, the exact-angle decomposition yields 3 non-Clifford rotation gates per controlled-$R_k$, but only a fraction trigger OFD attempts (the rest act on already-magic qubits within the MPS); the rank identity nevertheless provides the exact structural characterization of QFT's Class~II membership.

\subsection{Bond Dimension Bound Validation}
\label{sec:bound_tightness}

Figure~\ref{fig:bound_validation}(a) plots the final CAMPS bond dimension against GF(2) nullity for all 936 circuits.
Under the hybrid OFD+OBD strategy, all circuits satisfy $\chi \le 2^\nu$: OBD fallback keeps the bond dimension well below the theoretical ceiling.
However, the bound's tightness varies significantly across families (Fig.~\ref{fig:bound_validation}(b)).
Random circuits exhibit tightness ratios of $0.6$--$0.8$, meaning the bound is within a small constant factor; these circuits actually approach the worst case because OFD does not optimize the disentangler.

To validate the bound under pure OFD (no OBD fallback), we run 330 additional simulations at $N = 4$--$8$ across all 14 families.
The structural bound $\chi \le 2^{t-m}$ (Theorem~\ref{thm:bond_dim}, Claim~2), where $m$ is the number of greedy OFD successes, holds in 100\% of cases.
The stronger optimal bound $\chi \le 2^\nu$ holds in 96.7\% of cases; the remaining 3.3\% are circuits where greedy circuit-order OFD achieves $m < \mathrm{rank}(z)$ due to suboptimal qubit assignment (rank--$m$ deficit of 1--2, concentrated in random and QAOA circuits).
Theorem~\ref{thm:bond_dim} guarantees that an ordering achieving $m \ge \mathrm{rank}(z)$ \emph{exists} (Claim~1); greedy processing in circuit order does not always find it, but the structural bound $\chi \le 2^{t-m}$ holds for any ordering.

This looseness at small $N$ is a \emph{finite-size effect} rather than a limitation of the bound.
At small system sizes ($N = 4$--$16$), all circuits are classically simulable regardless of null space structure; bond dimensions remain manageable because the Hilbert space itself is small and OBD can efficiently explore the space of two-qubit Clifford disentanglers.
The bound's value is \emph{predictive}: it identifies which circuits will become exponentially expensive as $N$ grows.
At $N = 100$, a circuit with $\nu = 50$ has a bond dimension ceiling of $2^{50} \approx 10^{15}$. In contrast, when $\nu = 2$, the circuit remains trivially simulable, independent of OBD's ability to identify good disentanglers at small scale.
Moreover, even when OBD successfully keeps $\chi$ low at large $N$, the nullity directly controls the \emph{computational time}: each of the $\nu$ null-space non-Clifford gates requires OBD, so the total simulation time scales as $O(\nu \cdot t_{\mathrm{OBD}})$.
For QFT at $N = 32$ with $\nu/t = 0.98$ (Fig.~\ref{fig:qft_staircase}), nearly every non-Clifford gate triggers OBD. Consequently, the simulation is $10^4$--$10^5$ times slower than for a circuit of equal gate count in which OFD succeeds throughout, representing a computational barrier independent of the attainable bond dimension.
The small-$N$ benchmarks validate the \emph{correctness} of the bound; the scaling classification (Section~\ref{sec:classification}) establishes its \emph{relevance} at scale.

\subsection{Position-Dependent OFD Success}

\begin{figure}[t]
    \centering
    \includegraphics[width=\linewidth]{figure7.png}
    \caption{OFD success rate as a function of non-Clifford gate position within the circuit (normalized to $[0,1]$). Random circuits maintain high success throughout, while structured algorithms show position-dependent failure patterns. The declining success at later positions for random circuits reflects the gradual consumption of free qubits as the GF(2) row space fills.}
    \label{fig:ofd_position}
\end{figure}

Figure~\ref{fig:ofd_position} reveals position-dependent structure in OFD success.
Random circuits exhibit approximately $80$--$100\%$ success in early positions. The success rate declines toward the end of the circuit as free qubits are exhausted, consistent with the $\mathrm{GF}(2)$ row space approaching full rank at $\mathrm{rank} \approx N$.
Structured algorithms exhibit rapid initial decline: QFT, VQE, and Grover drop below 20\% success within the first 20\% of non-Clifford gates.
This reflects the fast saturation of rank: as confirmed by Proposition~\ref{prop:repeated_block}, the GF(2) rank reaches $N$ within the first block of non-Clifford gates, leaving all subsequent gates in the null space.

\subsection{Random Circuit Rank Saturation}

\begin{figure}[t]
    \centering
    \includegraphics[width=\linewidth]{figure8.png}
    \caption{Random Clifford+T rank saturation: GF(2) rank versus T-gate count for $N = 8$, $16$, and $32$.
    The rank increases linearly until $t \approx N$, then saturates sharply at $\mathrm{rank} = N$. This confirms Liu and Clark's result~\cite{Liu2024} and provides the empirical foundation for Class~I behavior: random circuits with $t = N$ achieve near-full rank with nullity $\approx 0.85$.}
    \label{fig:rank_saturation}
\end{figure}

Figure~\ref{fig:rank_saturation} shows the GF(2) rank as a function of T-gate count for random Clifford+T circuits at $N = 8$, $16$, and $32$.
The rank increases linearly for $t < N$, with each new T-gate adding approximately one independent row, then saturates sharply at $\mathrm{rank} = N$.
This saturation confirms Liu and Clark's analytical result~\cite{Liu2024} and the random matrix theory prediction: a random $t \times N$ binary matrix over $\mathbb{F}_2$ has expected rank $\min(t, N) - O(1)$, with the $O(1)$ correction being $\approx 0.85$ for square matrices.
The sharp saturation also validates our scaling classification: Class~I circuits have $\nu/t \to 0$ precisely because $\mathrm{rank}\approx N$ while $t = N$, giving $\nu \approx 0.85 = O(1)$.

\subsection{The Bond Dimension Paradox and Time-Space Decoupling}

Table~\ref{tab:ofd_rates} reveals a counterintuitive pattern: random circuits with high OFD success accumulate large bond dimensions ($\chi \approx 38$), while most structured algorithms maintain low bond dimensions ($\chi \le 6$), with the exception of QAOA ($\chi \approx 99$) whose low OFD rate ($39\%$) leaves the majority of rotation gates to the OBD fallback (see Remark~\ref{rem:qaoa} for the role of Clifford-angle rotations).
This apparent paradox resolves as follows.

The bound $\chi \le 2^\nu$ (Theorem~\ref{thm:bond_dim}) applies to pure OFD processing, where each failed disentangling doubles the bond dimension.
Our simulations employ OBD fallback (Section~\ref{sec:validation}), which successfully controls bond dimension for structured circuits at $N \le 16$: OBD searches over two-qubit Clifford operators to minimize the second R\'enyi entropy at each bond, actively reducing entanglement at every step.
For structured algorithms with regular entanglement patterns, OBD's optimization is highly effective, yielding near-optimal disentanglers that keep $\chi$ far below the $2^\nu$ ceiling.
Whether OBD maintains this effectiveness at larger $N$ remains an open question; the GF(2) nullity provides a worst-case bound independent of OBD performance.

The practical consequence is that high GF(2) nullity incurs a \emph{time cost} rather than a representational cost.
CAMPS simulations of structured algorithms remain compact (low $\chi$) but slow, as each of the $\nu$ OBD-requiring non-Clifford gates incurs the full OBD cost.
The $\mathrm{GF}(2)$ nullity thus predicts the \emph{time complexity} of CAMPS simulation, while the bond dimension controls the \emph{space complexity}, and these two costs are decoupled for structured circuits.
This motivates our GF(2)-guided improvements (Section~\ref{sec:practical}): knowing in advance which gates will fail OFD allows the simulation to proceed directly to OBD, saving the overhead of futile OFD attempts.


%% ============================================================
\section{Discussion}
\label{sec:discussion}
%% ============================================================

\textbf{Implications for CAMPS practitioners.}
Our scaling classification provides an immediate practical tool: before committing to an expensive CAMPS simulation, compute the GF(2) matrix and check the nullity ratio.
Class~I circuits will be fast; Class~II circuits will be slow regardless of implementation quality.
For Class~III circuits, the backbone sampling analysis (Section~\ref{sec:backbone}) provides quantitative guidance: if the Clifford backbone has full rank and $t = N$ T-gates are placed uniformly, the expected nullity ratio is $1/e \approx 0.368$, meaning roughly $63\%$ of T-gates will be OFD-amenable.
Circuits with partial backbones (e.g., Bell, Simon's) have higher nullity ratios, and QAOA requires separate analysis due to its qualitatively different structure.

\textbf{Connection to quantum advantage.}
The observation that structured quantum algorithms are precisely those with high GF(2) nullity has implications for quantum advantage claims.
Liu and Clark's result that random Clifford+T circuits with $t \le N$ are CAMPS-simulable~\cite{Liu2024} does \emph{not} extend to the structured circuits most relevant to practical quantum computing.
However, high nullity does not imply computational hardness. Rather, it indicates that \emph{OFD-based} simulation is slow, although alternative approaches such as OBD, STN~\cite{MasotLlima2024}, or Pauli propagation~\cite{PauliProp2019} may remain effective.
The bond dimension paradox demonstrates this directly: structured algorithms maintain low $\chi$ despite OFD failure.

\textbf{Relationship to nonstabilizerness measures.}
Frau \emph{et al.}~\cite{Frau2025} introduced a classification of many-body ground states based on how effectively local Clifford sweeps reduce entanglement entropy: fully locally-Clifford-disentanglable (fLCD) states can be perfectly disentangled, locally-Clifford-disentanglable (LCD) states exhibit entanglement reduction that improves with system size, and non-locally-Clifford-disentanglable (nLCD) states show only constant entanglement reduction.
Our three-class circuit taxonomy exhibits a suggestive parallel: Class~I circuits (where OFD succeeds at nearly all non-Clifford gates) loosely echo fLCD, Class~II circuits (where OFD fails universally) echo nLCD, and Class~III echoes LCD.
However, the two analogies apply at different levels of description. The classification of Frau \emph{et al.} concerns Clifford disentangling of quantum \emph{states}, as characterized by the mutual stabilizer R\'{e}nyi entropy. In contrast, our framework analyzes the $\mathrm{GF}(2)$ rank structure of \emph{circuits}.
Whether a formal connection exists between these two frameworks remains an open question.

\textbf{Beyond CAMPS.}
The GF(2) matrix encodes structural information about twisted Pauli strings that may be useful beyond the CAMPS framework.
In stabilizer tensor networks~\cite{MasotLlima2024}, each non-Clifford gate can increase the bond dimension $\chi$ of the tensor network of stabilizer amplitudes; we conjecture that the GF(2) null space structure may predict which non-Clifford gates cause this increase, though this connection has not been formally established.
Similarly, Pauli propagation methods~\cite{PauliProp2019} study Heisenberg-picture evolution, where non-Clifford gates generate branching of Pauli strings. Because the $\mathrm{GF}(2)$ row structure encodes the $X/Y$ content of the twisted Paulis governing each gate's interaction with the Clifford frame, we expect it to influence which gates contribute to branching, although this remains an open question.
Establishing formal connections between our GF(2) classification and the simulation costs of these alternative methods is an interesting direction for future work.

\textbf{Limitations.}
Our scaling analysis extends to $N = 64$ for most families, with QFT and VQE reaching $N = 256$.
While the GF(2) computation itself scales beyond $N = 256$, circuit generation for some families (e.g., Grover with many iterations) becomes the bottleneck.
The small-scale validation is limited to $N \le 16$ by CAMPS simulation cost.
Our formal results (Theorem~\ref{thm:bond_dim}, Propositions~\ref{prop:repeated_block}--\ref{prop:qft_rotation}, Theorem~\ref{thm:qft_rank}, Proposition~\ref{prop:backbone_bound}) are rigorous, with the lower bound of Theorem~\ref{thm:qft_rank} additionally supported by numerical verification at all $N$ from 4 to 256.
The backbone rank bound (Proposition~\ref{prop:backbone_bound}) applies to circuits with a static Clifford backbone; the birthday-effect corollary (Corollary~\ref{cor:birthday}) is verified empirically but relies on the uniform sampling assumption.
We benchmark only OFD/OBD disentangling strategies; the recently developed methods of Paviglianiti \emph{et al.}~\cite{Paviglianiti2024} and adaptive truncation strategies may modify the practical landscape.
The exact rank identity $\mathrm{rank} = N - 1$ (Theorem~\ref{thm:qft_rank}) is established for the standard controlled-rotation decomposition (Eq.~\eqref{eq:cr_decomp}) with exact angles; alternative decompositions (e.g., via gate synthesis into T-gates at finite precision $\varepsilon$) may yield different rank values, though the scaling class is unaffected since rank $= \Theta(N)$ in either case.

%% ============================================================
\section{Conclusion}
\label{sec:conclusion}
%% ============================================================

We have presented a GF(2) null space scaling classification of 14 quantum circuit families, revealing three distinct classes of CAMPS simulability: OFD-friendly (random circuits with $t = N$, $\nu/t \to 0$), OFD-hostile (QFT, VQE, Grover, $\nu/t \to 1$), and intermediate, whose backbone sampling structure predicts limiting nullity ratios via the classical birthday effect.
The classification is underpinned by formal results: a bond dimension bound $\chi \le 2^{\nu}$ (Theorem~\ref{thm:bond_dim}), a repeated-block nullity result explaining Class~II membership (Proposition~\ref{prop:repeated_block}), a Z-type obstruction characterization (Observation~\ref{prop:z_type}), a per-rotation rank bound (Proposition~\ref{prop:qft_rotation}), an exact rank identity $\mathrm{rank}(z_{\mathrm{QFT}}) = N - 1$ (Theorem~\ref{thm:qft_rank}) proved via the sequential structure of the QFT circuit, and a backbone rank bound for intermediate circuits (Proposition~\ref{prop:backbone_bound}) with a birthday-effect corollary predicting $\nu/t \to 1/e$ (Corollary~\ref{cor:birthday}).

Our key empirical finding is that GF(2) rank predicts CAMPS OFD success with Pearson $r = 0.95$ at the individual circuit level ($r = 0.98$ per-family), validated across 904 circuits and 14 families.
The structural bond dimension bound $\chi \le 2^{t-m}$ holds universally under pure OFD, with greedy circuit-order processing achieving the optimal bound $\chi \le 2^\nu$ in 96.7\% of tested circuits.
The GF(2) analysis runs in seconds even at $N = 256$, providing a fast and reliable predictor of CAMPS simulability \emph{before} expensive MPS simulation.
Practically, GF(2) pre-screening eliminates 95--99\% of futile OFD attempts for structured algorithms.

Looking forward, several directions remain open.
Our proof of the rank identity $\mathrm{rank}(z_{\mathrm{QFT}}) = N - 1$ is constructive: the lower bound identifies the specific gate (the control $R_z$ after the CNOT pair) in each iteration that produces $e_j$, with the mechanism fully determined by the CNOT column-update algebra.
Extending this structural analysis to additional circuit families would further develop the classification framework. Such work would require identifying the gate-level mechanisms that generate or inhibit independent $\mathrm{GF}(2)$ rows.
Extending the classification to Hamiltonian simulation circuits (Trotter, QDRIFT) would connect to the dominant use case of CAMPS for condensed matter physics~\cite{Qian2024,MelloTDVP2025}.
The GF(2) null space structure may also inform the development of new disentangling strategies specifically targeting Class~II circuits, potentially enabling efficient CAMPS simulation of structured quantum algorithms.

\section*{Code and Data Availability}

CAMPS.jl provides the first open-source implementation of the CAMPS framework. The associated repository contains the $\mathrm{GF}(2)$ analysis code, benchmarking scripts for all 14 circuit families, and data reproducing all figures, available at \url{https://github.com/riteshbhirud/CAMPS-OFD}.

\section*{Acknowledgments}
Acknowledgments....


%% ============================================================
%% BIBLIOGRAPHY
%% ============================================================

\bibliographystyle{plain}
\begin{thebibliography}{99}

\bibitem{Gottesman1998}
D.~Gottesman,
The Heisenberg representation of quantum computers,
in \textit{Proc.\ XXII International Colloquium on Group Theoretical Methods in Physics},
pp.~32--43 (1998).
arXiv:quant-ph/9807006.

\bibitem{Aaronson2004}
S.~Aaronson and D.~Gottesman,
Improved simulation of stabilizer circuits,
\textit{Phys.\ Rev.\ A} \textbf{70}, 052328 (2004).

\bibitem{Schollwock2011}
U.~Schollw\"{o}ck,
The density-matrix renormalization group in the age of matrix product states,
\textit{Ann.\ Phys.} \textbf{326}, 96--192 (2011).

\bibitem{White1992}
S.~R.~White,
Density matrix formulation for quantum renormalization groups,
\textit{Phys.\ Rev.\ Lett.} \textbf{69}, 2863 (1992).

\bibitem{Vidal2003}
G.~Vidal,
Efficient classical simulation of slightly entangled quantum computations,
\textit{Phys.\ Rev.\ Lett.} \textbf{91}, 147902 (2003).

\bibitem{Liu2024}
Z.~Liu and B.~K.~Clark,
Classical simulability of Clifford+T circuits with Clifford-augmented matrix product states,
arXiv:2412.17209 (2024).

\bibitem{Qian2024}
X.~Qian, J.~Huang, and M.~Qin,
Augmenting density matrix renormalization group with Clifford circuits,
\textit{Phys.\ Rev.\ Lett.} \textbf{133}, 190402 (2024).

\bibitem{MelloTDVP2025}
A.~F.~Mello, A.~Santini, G.~Lami, J.~De~Nardis, and M.~Collura,
Clifford dressed time-dependent variational principle,
\textit{Phys.\ Rev.\ Lett.} \textbf{134}, 150403 (2025).

\bibitem{Mello2024}
A.~F.~Mello, A.~Santini, and M.~Collura,
Hybrid stabilizer matrix product operator,
\textit{Phys.\ Rev.\ Lett.} \textbf{133}, 150604 (2024).

\bibitem{Qian2024b}
X.~Qian, J.~Huang, and M.~Qin,
Clifford circuits augmented time-dependent variational principle,
\textit{Phys.\ Rev.\ Lett.} \textbf{134}, 150404 (2025).

\bibitem{MasotLlima2024}
G.~Masot-Llima and A.~Garcia-Saez,
Stabilizer tensor networks: universal quantum simulator on a basis of stabilizer states,
\textit{Phys.\ Rev.\ Lett.} \textbf{133}, 230601 (2024).

\bibitem{PauliProp2019}
P.~Rall, D.~Liang, J.~Cook, and W.~Kretschmer,
Simulation of qubit quantum circuits via Pauli propagation,
\textit{Phys.\ Rev.\ A} \textbf{99}, 062337 (2019).

\bibitem{Zhang2025}
Y.~Zhang and Y.~Zhang,
Classical simulability of quantum circuits with shallow magic depth,
\textit{PRX Quantum} \textbf{6}, 010337 (2025).
arXiv:2409.13809.

\bibitem{Fux2025}
G.~E.~Fux, B.~B\'{e}ri, R.~Fazio, and E.~Tirrito,
Disentangling magic states with classically simulable quantum circuits,
\textit{Phys.\ Rev.\ Lett.} \textbf{135}, 260605 (2025).

\bibitem{Frau2025}
M.~Frau, P.~S.~Tarabunga, M.~Collura, E.~Tirrito, and M.~Dalmonte,
Stabilizer disentangling of conformal field theories,
\textit{SciPost Phys.} \textbf{18}, 165 (2025).

\bibitem{Huang2025}
J.~Huang, X.~Qian, and M.~Qin,
Nonstabilizerness entanglement entropy: a measure of hardness in the classical simulation of quantum many-body systems with tensor network states,
\textit{Phys.\ Rev.\ A} \textbf{112}, 012425 (2025).

\bibitem{Turkeshi2025}
X.~Turkeshi, A.~Dymarsky, and P.~Sierant,
Pauli spectrum and nonstabilizerness of typical quantum many-body states,
\textit{Phys.\ Rev.\ B} \textbf{111}, 054301 (2025).

\bibitem{Dehaene2003}
J.~Dehaene and B.~De~Moor,
The Clifford group, stabilizer states, and linear and quadratic operations over GF(2),
\textit{Phys.\ Rev.\ A} \textbf{68}, 042318 (2003).

\bibitem{Ross2016}
N.~J.~Ross and P.~Selinger,
Optimal ancilla-free Clifford+T approximation of $z$-rotations,
\textit{Quantum Inf.\ Comput.} \textbf{16}, 901--953 (2016).
arXiv:1403.2975.

\bibitem{Vandaele2025}
V.~Vandaele,
Lower T-count with faster algorithms,
\textit{Quantum} \textbf{9}, 1860 (2025).

\bibitem{Krastanov2022}
S.~Krastanov \textit{et al.},
QuantumClifford.jl: Clifford circuits, graph states, and stabilizer formalism tools,
Zenodo (2022).
\url{https://doi.org/10.5281/zenodo.6796187}.

\bibitem{Fishman2022}
M.~Fishman, S.~R.~White, and E.~M.~Stoudenmire,
The ITensor software library for tensor network calculations,
\textit{SciPost Phys.\ Codebases} \textbf{4} (2022).

\bibitem{Paviglianiti2024}
A.~Paviglianiti, G.~Lami, M.~Collura, and A.~Silva,
Estimating non-stabilizerness dynamics without simulating it,
arXiv:2405.06054 (2024).

\bibitem{BravyiGosset2016}
S.~Bravyi and D.~Gosset,
Improved classical simulation of quantum circuits dominated by Clifford gates,
\textit{Phys.\ Rev.\ Lett.} \textbf{116}, 250501 (2016).

\bibitem{Bravyi2019}
S.~Bravyi, D.~Browne, P.~Calpin, E.~Campbell, D.~Gosset, and M.~Howard,
Simulation of quantum circuits by low-rank stabilizer decompositions,
\textit{Quantum} \textbf{3}, 181 (2019).

\bibitem{Selinger2013}
P.~Selinger,
Quantum circuits of T-depth one,
\textit{Phys.\ Rev.\ A} \textbf{87}, 042302 (2013).

\bibitem{farhi2014quantum}
E.~Farhi, J.~Goldstone, and S.~Gutmann,
A quantum approximate optimization algorithm,
\textit{arXiv:1411.4028} (2014).

\end{thebibliography}

%% ============================================================
%% APPENDIX
%% ============================================================
\appendix

\section{GF(2) Matrix Construction from Stabilizer Tableaux}
\label{app:gf2_construction}

We provide explicit details for computing the GF(2) matrix $z$ from a stabilizer tableau, following the conventions of Aaronson and Gottesman~\cite{Aaronson2004}.

The stabilizer tableau for an $N$-qubit Clifford circuit $C$ is a $(2N) \times (2N+1)$ binary matrix.
Rows $1$ through $N$ encode the destabilizer generators, and rows $N+1$ through $2N$ encode the stabilizer generators.
The first $N$ columns store $x$-bits (presence of $X$), the next $N$ columns store $z$-bits (presence of $Z$), and the last column stores the phase.
For the identity circuit, row $i$ (destabilizer) encodes $X_i$ and row $N+i$ (stabilizer) encodes $Z_i$.

The Clifford gate update rules (applied to every row $i$) are:
\begin{align}
\text{Hadamard on $a$:}&\quad r_i \leftarrow r_i \oplus x_{ia}z_{ia};\; \text{swap}\; x_{ia} \leftrightarrow z_{ia} \\
\text{Phase ($S$) on $a$:}&\quad r_i \leftarrow r_i \oplus x_{ia}z_{ia};\; z_{ia} \leftarrow z_{ia} \oplus x_{ia} \\
\text{CNOT($a {\to} b$):}&\quad r_i \leftarrow r_i \oplus x_{ia}z_{ib}(x_{ib} \oplus z_{ia} \oplus 1) \notag\\
&\quad x_{ib} \leftarrow x_{ib} \oplus x_{ia};\; z_{ia} \leftarrow z_{ia} \oplus z_{ib}
\end{align}

To construct the GF(2) matrix: process the circuit gate by gate, maintaining a single stabilizer tableau.
For Clifford gates, apply the update rules above.
For each non-Clifford gate on qubit $q_k$ (non-Clifford gates are \emph{not} applied to the tableau):
\begin{enumerate}
\item Read row $N + q_k$ of the tableau (the stabilizer row for qubit $q_k$).
\item Extract the $x$-bits at positions $1$ through $N$: these form row $k$ of the GF(2) matrix $z$.
\item Do not modify the tableau---the non-Clifford gate's $Z$-component commutes through subsequent Clifford tracking.
\end{enumerate}

\noindent The key correctness argument: the cumulative Clifford circuit $C_k$ preceding the $k$-th non-Clifford gate includes all Clifford gates processed so far (including those between earlier non-Clifford gates).
The $x$-bits of row $N + q_k$ encode the image $C_k^\dagger Z_{q_k} C_k$, which is precisely the twisted Pauli $\bar{Z}^{[k]}$.
Skipping the non-Clifford gate in the tableau update is valid because $R_z(\theta)^\dagger Z R_z(\theta) = Z$ (any $R_z$ commutes with $Z$ up to global phase), so the presence or absence of non-Clifford gates in the Clifford accumulation does not affect the $x$-bits of subsequent twisted Paulis.

\section{GF(2) Gaussian Elimination}
\label{app:gf2_elimination}

The GF(2) rank is computed via Gaussian elimination over $\mathbb{F}_2$, where all arithmetic is modulo 2:

\begin{algorithmic}[1]
\Function{GF2Rank}{$M \in \mathbb{F}_2^{t \times N}$}
\State $r \gets 0$  \Comment{current rank}
\For{$c = 1, \ldots, N$}
    \State Find smallest $i > r$ with $M_{ic} = 1$
    \If{no such $i$} \textbf{continue} \EndIf
    \State $r \gets r + 1$; swap rows $r$ and $i$
    \For{each row $j \ne r$ with $M_{jc} = 1$}
        \State $M_{j,:} \gets M_{j,:} \oplus M_{r,:}$ \Comment{XOR row}
    \EndFor
\EndFor
\State \Return $r$
\EndFunction
\end{algorithmic}

\noindent The cost is $O(\min(t,N)^2 \cdot \max(t,N))$ bit operations.
For most circuits in this study ($N \le 64$, $t \le 7{,}000$), this completes in milliseconds.
For the largest configurations ($N = 256$, $t \approx 80{,}000$), we use an online RREF variant that maintains an $N \times N$ matrix in reduced row echelon form, inserting rows one at a time in $O(Nr)$ bit operations per row where $r \le N$ is the current rank, giving $O(N^2 t)$ total time and $O(N^2)$ memory.

\end{document}